\chapter[Gravitational Endpoint Transfer of the Active FCC Vacuum]{Gravitational Endpoint Transfer\\[0.5ex]
\Large Gravitational Constraint, Coordinate Reduction, Referred Inertia, and Dispersion}
\label{app:gravitational_endpoint_transfer}

\noindent\textbf{Subject:}
Gravity-specific constrained-coordinate reduction of the active FCC Soft Floor and derivation of its referred inertia and long-wavelength propagation law.

\section{Scope and Status of the Construction}
\label{sec:grav_endpoint_scope}

\hyperref[app:fcc_microphysics]{Appendix~\ref*{app:fcc_microphysics}} develops the gravity-independent FCC contact geometry, objective bond kinematics, parity-even contact action, primitive-cell assembly, and long-wavelength acoustic tensor of the active Soft Floor. The present appendix applies the gravitational loading constraint to that baseline mechanical carrier.

The construction is restricted to the active FCC phase,
\begin{equation}
    0 \leq \alpha_{\mathrm{G}} < 1,
\end{equation}
and to a gravitational background that is quasistatic during one local transfer event,
\begin{equation}
    \dot{\alpha}_{\mathrm{G}} \approx 0.
\end{equation}

The analysis develops the sequence
\begin{equation}
\begin{array}{c}
    \text{gravitational capacity reduction} \\
    \longrightarrow
    \text{work-conjugate endpoint map} \\
    \longrightarrow
    \text{serial source-target constraint} \\
    \longrightarrow
    \text{latent-to-completed coordinate reduction} \\
    \longrightarrow
    \text{referred inertia} \\
    \longrightarrow
    \text{constrained FCC dispersion}.
\end{array}
\end{equation}

The single-endpoint map is a reduced work-conjugate constitutive construction. It is not claimed to have been derived from a complete parent potential containing all node, contact, condensate, phase, and Tension-Column degrees of freedom. The present appendix instead shows that, once the endpoint map is defined through the active-capacity balance, its inertial and dispersive consequences follow from a conservative constrained action.

The principal statuses are:

\begin{table}[h!]
\centering
\renewcommand{\arraystretch}{1.3}
\begin{tabular}{p{0.45\linewidth} p{0.5\linewidth}}
\hline
\textbf{Result} & \textbf{Status in this appendix} \\
\hline
$\tau_{\mathrm{av}}=\tau_{\mathrm{D}}-\tau_{\mathrm{L}}$ &
active-capacity subtraction within the reduced constitutive model
\\
$\mathcal{C}_{i}=1-\alpha_{\mathrm{G}}^{(i)}$ &
derived from the work-conjugate endpoint map
\\
$\mathcal{C}_{ij}=\mathcal{C}_{i}\mathcal{C}_{j}$ &
derived by serial source release and target acceptance
\\
$u=\mathcal{C}y$ &
holonomic latent-to-completed coordinate constraint
\\
$\mathbf{M}_{\mathrm{cell}}^{(\mathrm{eff})}=\mathbf{T}_{\mathrm{cell}}^{-\mathsf{T}}\mathbf{M}_{\mathrm{cell}}^{(0)}\mathbf{T}_{\mathrm{cell}}^{-1}$ &
mass-matrix congruence within the stated quasistatic holonomic reduction
\\
$\rho_{\mathrm{eff}}=\rho_{\mathrm{eff}}^{(0)}/\mathcal{C}^{2}$ &
referred inertia in the coherent translational sector
\\
$c_{\mathrm{P}}=c_{\mathrm{E}}\mathcal{C}$ &
derived from the constrained directional and FCC acoustic equations
\\
full parent microscopic origin of $\mathcal{C}_{i}$ &
not yet completed
\\
\hline
\end{tabular}
\end{table}

\section{Gravitational Loading and Available Active Capacity}
\label{sec:grav_active_capacity}

For endpoint $i$, define the dimensionless gravitational loading fraction by
\begin{equation}
    \alpha_{\mathrm{G}}^{(i)} = \frac{\tau_{\mathrm{L}}^{(i)}}{\tau_{\mathrm{D}}},
\end{equation}
where $\tau_{\mathrm{L}}^{(i)}$ is the local mass-induced loading and $\tau_{\mathrm{D}}$ is the baseline active surface-tension capacity.

The active capacity remaining for additional local transfer is
\begin{equation}
    \tau_{\mathrm{av}}^{(i)} = \tau_{\mathrm{D}}-\tau_{\mathrm{L}}^{(i)}.
\end{equation}

Substituting the definition of $\alpha_{\mathrm{G}}^{(i)}$ gives
\begin{equation}
    \tau_{\mathrm{av}}^{(i)} = \tau_{\mathrm{D}} \left( 1-\alpha_{\mathrm{G}}^{(i)} \right).
\end{equation}

Define the single-endpoint residual-capacity factor
\begin{equation}
    \mathcal{C}_{i} \equiv 1-\alpha_{\mathrm{G}}^{(i)}.
    \label{eq:single_endpoint_capacity}
\end{equation}

Then
\begin{equation}
    \tau_{\mathrm{av}}^{(i)} = \tau_{\mathrm{D}}\mathcal{C}_{i}.
\end{equation}

Within the active phase,
\begin{equation}
    0 < \mathcal{C}_{i} \leq 1.
\end{equation}

The limits are
\begin{equation}
    \alpha_{\mathrm{G}}^{(i)} = 0 \quad\Longrightarrow\quad \mathcal{C}_{i} = 1,
\end{equation}
and
\begin{equation}
    \alpha_{\mathrm{G}}^{(i)} \rightarrow 1^{-} \quad\Longrightarrow\quad \mathcal{C}_{i} \rightarrow 0^{+}.
\end{equation}

The linear dependence in Equation~\eqref{eq:single_endpoint_capacity} is therefore not inferred from the limiting values alone. It follows from the adopted linear subtraction of the local gravitational loading from the available active tension capacity.

\section{Work-Conjugate Single-Endpoint Coordinate Map}
\label{sec:grav_work_conjugate_endpoint}

Let $y_{i}$ denote the latent node--condensate reconfiguration presented to endpoint $i$, and let $u_{i}$ denote the completed displacement accepted by the active FCC contact state.

Introduce the dimensionless normalized coordinates
\begin{equation}
    \widehat{y}_{i} \equiv \frac{y_{i}}{d_{\mathrm{nn}}}, \qquad \widehat{u}_{i} \equiv \frac{u_{i}}{d_{\mathrm{nn}}},
\end{equation}
where
\begin{equation}
    d_{\mathrm{nn}} = 2r_{\mathrm{D}}
\end{equation}
is the nearest-neighbor FCC separation.

The coordinate $\widehat{y}_{i}$ measures the normalized reconfiguration presented to the endpoint, while $\widehat{u}_{i}$ measures the normalized deformation transmitted into the active FCC contact state.

The reduced endpoint closure is obtained by equating the incremental work density represented in the two coordinate descriptions:
\begin{equation}
    \tau_{\mathrm{D}}\,\delta\widehat{u}_{i} = \tau_{\mathrm{av}}^{(i)}\,\delta\widehat{y}_{i}.
    \label{eq:endpoint_work_conjugacy}
\end{equation}

Using
\begin{equation}
    \tau_{\mathrm{av}}^{(i)} = \tau_{\mathrm{D}}\mathcal{C}_{i},
\end{equation}
Equation~\eqref{eq:endpoint_work_conjugacy} becomes
\begin{equation}
    \tau_{\mathrm{D}}\,\delta\widehat{u}_{i} = \tau_{\mathrm{D}}\mathcal{C}_{i}\,\delta\widehat{y}_{i}.
\end{equation}

Therefore,
\begin{equation}
    \delta\widehat{u}_{i} = \mathcal{C}_{i}\,\delta\widehat{y}_{i}.
\end{equation}

For a quasistatic gravitational background, $\mathcal{C}_{i}$ is constant during the local transfer event. Integrating from the common undeformed reference state gives
\begin{equation}
    \widehat{u}_{i} = \mathcal{C}_{i}\widehat{y}_{i}.
\end{equation}

Multiplying by $d_{\mathrm{nn}}$ gives the single-endpoint coordinate map
\begin{equation}
    u_{i} = \mathcal{C}_{i}y_{i}.
    \label{eq:single_endpoint_coordinate_map}
\end{equation}

Thus $\mathcal{C}_{i}$ is the fraction of the latent reconfiguration admitted into the completed active-FCC coordinate.

\section{Serial Source--Target Composition}
\label{sec:grav_serial_composition}

Consider a completed nearest-neighbor transfer from source endpoint $i$ to target endpoint $j$.

Let $y_{ij}$ denote the latent transfer amplitude presented to the source endpoint. The source releases only the fraction permitted by its residual active capacity:
\begin{equation}
    y_{ij}^{(\mathrm{rel})} = \mathcal{C}_{i}y_{ij}.
\end{equation}

The target endpoint acts on the amplitude already released by the source:
\begin{equation}
    u_{ij} = \mathcal{C}_{j}y_{ij}^{(\mathrm{rel})}.
\end{equation}

Substitution gives
\begin{equation}
    u_{ij} = \mathcal{C}_{i}\mathcal{C}_{j}y_{ij}.
\end{equation}

Define the completed two-endpoint transfer factor by
\begin{equation}
    \mathcal{C}_{ij} \equiv \mathcal{C}_{i}\mathcal{C}_{j}.
    \label{eq:two_endpoint_factor}
\end{equation}

Using Equation~\eqref{eq:single_endpoint_capacity},
\begin{equation}
    \mathcal{C}_{ij} = \left( 1-\alpha_{\mathrm{G}}^{(i)} \right) \left( 1-\alpha_{\mathrm{G}}^{(j)} \right).
    \label{eq:gravity_endpoint_product}
\end{equation}

The completed coordinate therefore satisfies
\begin{equation}
    u_{ij} = \mathcal{C}_{ij}y_{ij}.
    \label{eq:two_endpoint_coordinate_map}
\end{equation}

The product is not a probability rule. It is the composition of two serial coordinate constraints: source release followed by target acceptance and closure.

For a gravitational loading that varies negligibly across one FCC connection,
\begin{equation}
    \alpha_{\mathrm{G}}^{(i)} \approx \alpha_{\mathrm{G}}^{(j)} \approx \alpha_{\mathrm{G}},
\end{equation}
and the local two-endpoint factor reduces to
\begin{equation}
    \mathcal{C} = \left( 1-\alpha_{\mathrm{G}} \right)^2.
    \label{eq:uniform_two_endpoint_factor}
\end{equation}

\section{Constrained Directional Action}
\label{sec:grav_constrained_directional_action}

\hyperref[app:fcc_microphysics]{Appendix~\ref*{app:fcc_microphysics}} establishes the baseline transverse nearest-neighbor action of the unstrained FCC carrier. The gravitational specialization distinguishes the latent reconfiguration coordinate $y_{n}$ from the completed active-FCC displacement $u_{n}$.

For a locally uniform, quasistatic gravitational loading, impose the holonomic constraint
\begin{equation}
    u_{n} = \mathcal{C}y_{n}.
    \label{eq:directional_holonomic_constraint}
\end{equation}

A constrained directional Lagrangian per unit participating interface area is
\begin{equation}
\begin{aligned}
    \mathcal{L}_{\perp} = \sum_{n} \Bigg[ & \frac{1}{2}\mu_{\mathrm{D}}\dot{y}_{n}^2 - \frac{\tau_{\mathrm{D}}}{2d_{\mathrm{nn}}^2} \left( u_{n+1}-u_{n} \right)^2 \\
    & + \lambda_{n} \left( u_{n}-\mathcal{C}y_{n} \right) \Bigg].
    \label{eq:constrained_directional_lagrangian}
\end{aligned}
\end{equation}

Variation with respect to $\lambda_{n}$ gives Equation~\eqref{eq:directional_holonomic_constraint}. Since
\begin{equation}
    y_{n} = \frac{u_{n}}{\mathcal{C}},
\end{equation}
and
\begin{equation}
    \dot{\mathcal{C}} \approx 0,
\end{equation}
one has
\begin{equation}
    \dot{y}_{n} = \frac{\dot{u}_{n}}{\mathcal{C}}.
\end{equation}

Eliminating the latent coordinate gives the reduced completed-coordinate Lagrangian
\begin{equation}
    \mathcal{L}_{\perp}^{(\mathrm{red})} = \sum_{n} \left[ \frac{1}{2} \frac{\mu_{\mathrm{D}}}{\mathcal{C}^2} \dot{u}_{n}^2 - \frac{\tau_{\mathrm{D}}}{2d_{\mathrm{nn}}^2} \left( u_{n+1}-u_{n} \right)^2 \right].
    \label{eq:reduced_directional_lagrangian}
\end{equation}

The active-phase tangent stiffness remains the baseline FCC stiffness at the retained constitutive order. The gravitational constraint instead modifies the kinetic coefficient referred to the completed coordinate.

\section{Referred Surface Inertia}
\label{sec:grav_referred_surface_inertia}

Comparison of Equation~\eqref{eq:reduced_directional_lagrangian} with the standard kinetic form identifies the effective completed-coordinate surface inertia:
\begin{equation}
    \mu_{\mathrm{eff}} = \frac{\mu_{\mathrm{D}}}{\mathcal{C}^2}.
    \label{eq:surface_referred_inertia_general}
\end{equation}

For the locally uniform two-endpoint factor,
\begin{equation}
    \mathcal{C} = \left( 1-\alpha_{\mathrm{G}} \right)^2,
\end{equation}
the referred surface inertia becomes
\begin{equation}
    \mu_{\mathrm{eff}} = \frac{\mu_{\mathrm{D}}} {\left( 1-\alpha_{\mathrm{G}} \right)^4}.
    \label{eq:surface_referred_inertia_alpha}
\end{equation}

The inverse fourth power is therefore not introduced by algebraically solving an already assumed wave-speed relation. It follows directly from the constrained kinetic-coordinate transformation:
\begin{equation}
    u = \mathcal{C}y \quad\Longrightarrow\quad \dot{y} = \frac{\dot{u}}{\mathcal{C}} \quad\Longrightarrow\quad \mu_{\mathrm{eff}} = \frac{\mu_{\mathrm{D}}}{\mathcal{C}^2}.
\end{equation}

The coefficient $\mu_{\mathrm{eff}}$ is the referred surface inertia of the directional contact representation, whereas $\rho_{\mathrm{eff}}$ denotes the corresponding volumetric inertia coefficient at a specified continuum scale. The two quantities are related through the appropriate surface-to-volume projection and are not interchangeable.

\section{Primitive-Cell Transfer Operator}
\label{sec:grav_cell_transfer_operator}

The scalar endpoint constraint must be embedded into the complete primitive-cell coordinate space.

Let $\mathbf{y}_{\mathrm{cell}}$ denote the latent generalized primitive-cell coordinate, and let $\mathbf{q}_{\mathrm{cell}}$ denote the corresponding completed active-FCC coordinate.

Let $\mathbb{P}_{\mathrm{tr}}$ be the projector onto the coherent translational transfer sector. It satisfies
\begin{equation}
    \mathbb{P}_{\mathrm{tr}}^2 = \mathbb{P}_{\mathrm{tr}}, \qquad \mathbb{P}_{\mathrm{tr}}^{\mathsf{T}} = \mathbb{P}_{\mathrm{tr}}.
\end{equation}

The local gravitational transfer operator is
\begin{equation}
    \mathbf{T}_{\mathrm{cell}} = \mathbf{I} + \left( \mathcal{C}-1 \right) \mathbb{P}_{\mathrm{tr}}.
    \label{eq:cell_transfer_operator}
\end{equation}

The completed and latent coordinates are related by
\begin{equation}
    \mathbf{q}_{\mathrm{cell}} = \mathbf{T}_{\mathrm{cell}} \mathbf{y}_{\mathrm{cell}}.
    \label{eq:cell_coordinate_map}
\end{equation}

The operator acts as multiplication by $\mathcal{C}$ on the coherent translational sector and as the identity on the complementary retained sectors:
\begin{equation}
    \mathbf{T}_{\mathrm{cell}} \mathbb{P}_{\mathrm{tr}} = \mathcal{C}\mathbb{P}_{\mathrm{tr}},
\end{equation}
and
\begin{equation}
    \mathbf{T}_{\mathrm{cell}} \left( \mathbf{I}-\mathbb{P}_{\mathrm{tr}} \right) = \mathbf{I}-\mathbb{P}_{\mathrm{tr}}.
\end{equation}

For
\begin{equation}
    \mathcal{C} > 0,
\end{equation}
the inverse operator is
\begin{equation}
    \mathbf{T}_{\mathrm{cell}}^{-1} = \mathbf{I} + \left( \mathcal{C}^{-1}-1 \right) \mathbb{P}_{\mathrm{tr}}.
    \label{eq:cell_transfer_inverse}
\end{equation}

This follows directly from the idempotence of $\mathbb{P}_{\mathrm{tr}}$.

\section{Primitive-Cell Mass-Matrix Congruence}
\label{sec:grav_mass_matrix_congruence}

Let $\mathbf{M}_{\mathrm{cell}}^{(0)}$ be the gravity-neutral primitive-cell mass matrix constructed in \hyperref[app:fcc_microphysics]{Appendix~\ref*{app:fcc_microphysics}}. The latent-coordinate kinetic energy is
\begin{equation}
    T_{\mathrm{cell}} = \frac{1}{2} \dot{\mathbf{y}}_{\mathrm{cell}}^{\mathsf{T}} \mathbf{M}_{\mathrm{cell}}^{(0)} \dot{\mathbf{y}}_{\mathrm{cell}}.
\end{equation}

Using
\begin{equation}
    \mathbf{y}_{\mathrm{cell}} = \mathbf{T}_{\mathrm{cell}}^{-1} \mathbf{q}_{\mathrm{cell}},
\end{equation}
and the quasistatic condition
\begin{equation}
    \dot{\mathbf{T}}_{\mathrm{cell}} \approx \mathbf{0},
\end{equation}
the kinetic energy becomes
\begin{equation}
    T_{\mathrm{cell}} = \frac{1}{2} \dot{\mathbf{q}}_{\mathrm{cell}}^{\mathsf{T}} \mathbf{T}_{\mathrm{cell}}^{-\mathsf{T}} \mathbf{M}_{\mathrm{cell}}^{(0)} \mathbf{T}_{\mathrm{cell}}^{-1} \dot{\mathbf{q}}_{\mathrm{cell}}.
\end{equation}

Therefore the completed-coordinate mass matrix is
\begin{equation}
    \mathbf{M}_{\mathrm{cell}}^{(\mathrm{eff})} = \mathbf{T}_{\mathrm{cell}}^{-\mathsf{T}} \mathbf{M}_{\mathrm{cell}}^{(0)} \mathbf{T}_{\mathrm{cell}}^{-1}.
    \label{eq:effective_cell_mass_matrix}
\end{equation}

Equation~\eqref{eq:effective_cell_mass_matrix} is the mass-matrix congruence associated with the stated quasistatic holonomic transfer constraint.

Under the primitive-cell bookkeeping of \hyperref[app:fcc_microphysics]{Appendix~\ref*{app:fcc_microphysics}}, the coherent translational mass block is isotropic and decoupled from the retained complementary coordinates at leading order. In that sector,
\begin{equation}
    \mathbb{P}_{\mathrm{tr}} \mathbf{M}_{\mathrm{cell}}^{(\mathrm{eff})} \mathbb{P}_{\mathrm{tr}} = \frac{1}{\mathcal{C}^2} \mathbb{P}_{\mathrm{tr}} \mathbf{M}_{\mathrm{cell}}^{(0)} \mathbb{P}_{\mathrm{tr}}.
    \label{eq:projected_mass_scaling}
\end{equation}

Thus the cell-level operator reproduces the scalar referred-inertia result while preserving the unmodified complementary sectors at the retained order.

\section{Constrained FCC Dynamic Matrix}
\label{sec:grav_constrained_dynamic_matrix}

\hyperref[app:fcc_microphysics]{Appendix~\ref*{app:fcc_microphysics}} defines the gravity-neutral FCC Bloch stiffness
\begin{equation}
    \mathbf{K}_{\mathrm{FCC}}^{(0)}(\mathbf{k}) = A_{\mathrm{ct}} \sum_{a=1}^{6} \mathbf{B}_{a}^{\dagger}(\mathbf{k}) \mathbf{H}_{\mathrm{c}} \mathbf{B}_{a}(\mathbf{k}).
\end{equation}

At the retained active-phase tangent order, the gravitational coordinate reduction acts on the kinetic operator while the baseline completed-coordinate contact stiffness remains
\begin{equation}
    \mathbf{K}_{\mathrm{FCC}} \left( \mathbf{k}, \alpha_{\mathrm{G}} \right) = \mathbf{K}_{\mathrm{FCC}}^{(0)}(\mathbf{k}).
    \label{eq:baseline_stiffness_retained}
\end{equation}

The gravity-constrained generalized eigenvalue problem is therefore
\begin{equation}
    \det \left[ \mathbf{K}_{\mathrm{FCC}}^{(0)}(\mathbf{k}) - \omega^2 \mathbf{M}_{\mathrm{cell}}^{(\mathrm{eff})} \right] = 0.
    \label{eq:constrained_fcc_eigenproblem}
\end{equation}

The gravitational dependence enters through the congruence-transformed mass matrix rather than through an independent softening of the baseline tangent contact stiffness.

This avoids counting the same gravitational reduction simultaneously in both the stiffness and inertia sectors.

\section{Directional Lattice Dispersion}
\label{sec:grav_directional_dispersion}

Variation of the reduced directional Lagrangian in Equation~\eqref{eq:reduced_directional_lagrangian} gives
\begin{equation}
    \frac{\mu_{\mathrm{D}}}{\mathcal{C}^2} \ddot{u}_{n} = \frac{\tau_{\mathrm{D}}}{d_{\mathrm{nn}}^2} \left( u_{n+1}-2u_{n}+u_{n-1} \right).
    \label{eq:constrained_directional_equation}
\end{equation}

For a plane wave,
\begin{equation}
    u_{n}(t) = u_{0} \exp \left[ i \left( knd_{\mathrm{nn}}-\omega t \right) \right],
\end{equation}
one obtains
\begin{equation}
    u_{n+1}-2u_{n}+u_{n-1} = -4u_{n} \sin^2 \left( \frac{kd_{\mathrm{nn}}}{2} \right).
\end{equation}

The constrained directional dispersion relation is therefore
\begin{equation}
    \omega^2 = \frac{ 4\tau_{\mathrm{D}}\mathcal{C}^2 }{ \mu_{\mathrm{D}}d_{\mathrm{nn}}^2 } \sin^2 \left( \frac{kd_{\mathrm{nn}}}{2} \right).
\end{equation}

Using
\begin{equation}
    c_{\mathrm{E}}^2 = \frac{\tau_{\mathrm{D}}}{\mu_{\mathrm{D}}},
\end{equation}
the dispersion becomes
\begin{equation}
    \omega^2 \left( k, \alpha_{\mathrm{G}} \right) = \frac{ 4c_{\mathrm{E}}^2\mathcal{C}^2 }{ d_{\mathrm{nn}}^2 } \sin^2 \left( \frac{kd_{\mathrm{nn}}}{2} \right).
    \label{eq:constrained_directional_dispersion}
\end{equation}

In the long-wavelength limit,
\begin{equation}
    kd_{\mathrm{nn}} \ll 1,
\end{equation}
Equation~\eqref{eq:constrained_directional_dispersion} gives
\begin{equation}
    \omega^2 = c_{\mathrm{E}}^2\mathcal{C}^2k^2 + \mathcal{O} \left( k^4d_{\mathrm{nn}}^2 \right).
\end{equation}

Therefore,
\begin{equation}
    \lim_{k\rightarrow0} \frac{\omega}{k} = c_{\mathrm{E}}\mathcal{C}.
    \label{eq:directional_longwave_speed}
\end{equation}

Using Equation~\eqref{eq:uniform_two_endpoint_factor},
\begin{equation}
    c_{\mathrm{P}} = c_{\mathrm{E}} \left( 1-\alpha_{\mathrm{G}} \right)^2.
    \label{eq:gravity_propagation_law}
\end{equation}

Thus the propagation law follows through the mechanical sequence
\begin{equation}
    \text{work-conjugate endpoint map} \longrightarrow \mathcal{C} \longrightarrow \text{referred inertia} \longrightarrow \text{discrete dispersion} \longrightarrow c_{\mathrm{P}}.
\end{equation}

\section{Full FCC Long-Wavelength Reduction}
\label{sec:grav_full_fcc_reduction}

The directional derivation must agree with the three-dimensional FCC acoustic tensor developed in \hyperref[app:fcc_microphysics]{Appendix~\ref*{app:fcc_microphysics}}.

Let $K_{\mathrm{sh}}^{(0)}$ and $\rho_{\mathrm{eff}}^{(0)}$ denote the matched baseline shear coefficient and coherent translational inertia coefficient at one chosen descriptive scale. They satisfy
\begin{equation}
    \frac{K_{\mathrm{sh}}^{(0)}}{\rho_{\mathrm{eff}}^{(0)}} = c_{\mathrm{E}}^2.
    \label{eq:baseline_scale_matched_ratio}
\end{equation}

At the primitive-cell microscopic scale presently derived in \hyperref[app:fcc_microphysics]{Appendix~\ref*{app:fcc_microphysics}}, these scale-neutral coefficients are identified with
\begin{equation}
    \rho_{\mathrm{eff}}^{(0)} = \rho_{\mathrm{sh}}^{(\mathrm{micro},0)}, \qquad K_{\mathrm{sh}}^{(0)} = K_{\mathrm{sh}}^{(\mathrm{micro},0)}.
    \label{eq:microscopic_baseline_identification}
\end{equation}

Here $\rho_{\mathrm{sh}}^{(\mathrm{micro},0)}$ is the primitive-cell FCC scaffold inertia normalization derived in Appendix D. The identification in Equation~\eqref{eq:microscopic_baseline_identification} does not introduce an independently additive condensate kinetic density. Coherent condensate entrainment fixes the relative kinematics, while any additional coherent condensate inertia would require a separately derived coupled two-field kinetic action.

The gravitationally referred density in the coherent transverse sector is
\begin{equation}
    \rho_{\mathrm{eff}} \left( \alpha_{\mathrm{G}} \right) = \frac{ \rho_{\mathrm{eff}}^{(0)} }{ \mathcal{C}^2 }.
    \label{eq:referred_density_general}
\end{equation}

At the retained tangent order,
\begin{equation}
    K_{\mathrm{sh}} \left( \alpha_{\mathrm{G}} \right) = K_{\mathrm{sh}}^{(0)}.
\end{equation}

The long-wavelength transverse acoustic equation is therefore
\begin{equation}
    \omega_{\mathrm{T}}^2 = \frac{ K_{\mathrm{sh}}^{(0)} }{ \rho_{\mathrm{eff}} \left( \alpha_{\mathrm{G}} \right) } k^2 + \mathcal{O} \left( k^4d_{\mathrm{nn}}^2 \right).
\end{equation}

Substituting Equation~\eqref{eq:referred_density_general} gives
\begin{equation}
    \omega_{\mathrm{T}}^2 = \frac{ K_{\mathrm{sh}}^{(0)} }{ \rho_{\mathrm{eff}}^{(0)} } \mathcal{C}^2k^2 + \mathcal{O} \left( k^4d_{\mathrm{nn}}^2 \right).
\end{equation}

Using Equation~\eqref{eq:baseline_scale_matched_ratio},
\begin{equation}
    \omega_{\mathrm{T}}^2 = c_{\mathrm{E}}^2 \mathcal{C}^2k^2 + \mathcal{O} \left( k^4d_{\mathrm{nn}}^2 \right).
    \label{eq:fcc_constrained_transverse_dispersion}
\end{equation}

For locally uniform gravitational loading,
\begin{equation}
    \omega_{\mathrm{T}}^2 = c_{\mathrm{E}}^2 \left( 1-\alpha_{\mathrm{G}} \right)^4 k^2 + \mathcal{O} \left( k^4d_{\mathrm{nn}}^2 \right).
    \label{eq:fcc_alpha_dispersion}
\end{equation}

Consequently,
\begin{equation}
    c_{\mathrm{P}} = c_{\mathrm{E}} \left( 1-\alpha_{\mathrm{G}} \right)^2.
\end{equation}

The same gravitational factor therefore follows from both the constrained directional sequence and the full FCC long-wavelength acoustic tensor.

The baseline coefficients in this section must be taken at the same descriptive scale. Primitive-cell microscopic coefficients must not be mixed with independently inferred macroscopic continuum coefficients.

\section{Exact Nonuniform Endpoint Correction}
\label{sec:grav_nonuniform_endpoint}

For a selected FCC connection, define the local mean loading by
\begin{equation}
    \alpha_{\mathrm{G}} = \frac{ \alpha_{\mathrm{G}}^{(i)} + \alpha_{\mathrm{G}}^{(j)} }{2},
\end{equation}
and the endpoint difference by
\begin{equation}
    \delta\alpha_{\mathrm{G}} = \alpha_{\mathrm{G}}^{(j)} - \alpha_{\mathrm{G}}^{(i)}.
\end{equation}

Then
\begin{equation}
    \alpha_{\mathrm{G}}^{(i)} = \alpha_{\mathrm{G}} - \frac{\delta\alpha_{\mathrm{G}}}{2}, \qquad \alpha_{\mathrm{G}}^{(j)} = \alpha_{\mathrm{G}} + \frac{\delta\alpha_{\mathrm{G}}}{2}.
\end{equation}

Substitution into Equation~\eqref{eq:gravity_endpoint_product} gives
\begin{equation}
\begin{aligned}
    \mathcal{C}_{ij} &= \left( 1-\alpha_{\mathrm{G}} + \frac{\delta\alpha_{\mathrm{G}}}{2} \right) \left( 1-\alpha_{\mathrm{G}} - \frac{\delta\alpha_{\mathrm{G}}}{2} \right) \\
    &= \left( 1-\alpha_{\mathrm{G}} \right)^2 - \frac{ \left( \delta\alpha_{\mathrm{G}} \right)^2 }{4}.
\end{aligned}
\label{eq:exact_nonuniform_factor}
\end{equation}

The first-order endpoint-gradient contribution cancels exactly.

The bond-resolved propagation speed is
\begin{equation}
    c_{\mathrm{P}}^{(i\rightarrow j)} = c_{\mathrm{E}} \left[ \left( 1-\alpha_{\mathrm{G}} \right)^2 - \frac{ \left( \delta\alpha_{\mathrm{G}} \right)^2 }{4} \right].
    \label{eq:nonuniform_bond_speed}
\end{equation}

The locally uniform continuum approximation is controlled by
\begin{equation}
    \frac{ \left( \delta\alpha_{\mathrm{G}} \right)^2 }{4} \ll \left( 1-\alpha_{\mathrm{G}} \right)^2,
\end{equation}
or equivalently,
\begin{equation}
    \left| \delta\alpha_{\mathrm{G}} \right| \ll 2 \left| 1-\alpha_{\mathrm{G}} \right|.
    \label{eq:continuum_validity_condition}
\end{equation}

For a smooth spherical exterior field,
\begin{equation}
    \alpha_{\mathrm{G}}(R) = \frac{GM}{Rc_{\mathrm{E}}^2},
\end{equation}
the relative change across one radial nearest-neighbor separation is
\begin{equation}
    \frac{ \left| \delta\alpha_{\mathrm{G}} \right| }{ \alpha_{\mathrm{G}} } \sim \frac{2r_{\mathrm{D}}}{R}.
\end{equation}

The continuum correction is therefore negligible throughout ordinary macroscopic exterior regions. It cannot, however, be neglected automatically near the saturation boundary, where $1-\alpha_{\mathrm{G}}$ becomes small, or across a sharp FCC-to-SC phase interface.

\section{Refraction and Effective Optical Index}
\label{sec:grav_refraction_appendix}

The local propagation index of the coherent transverse branch is
\begin{equation}
    n_{\tau} = \frac{c_{\mathrm{E}}}{c_{\mathrm{P}}}.
\end{equation}

For locally uniform gravitational loading,
\begin{equation}
    n_{\tau} = \frac{1}{ \left( 1-\alpha_{\mathrm{G}} \right)^2 }.
    \label{eq:gravity_refraction_index_appendix}
\end{equation}

In the weak-field limit,
\begin{equation}
    \alpha_{\mathrm{G}} \ll 1,
\end{equation}
the expansion is
\begin{equation}
    n_{\tau} = 1 + 2\alpha_{\mathrm{G}} + 3\alpha_{\mathrm{G}}^2 + \mathcal{O} \left( \alpha_{\mathrm{G}}^3 \right).
    \label{eq:gravity_index_expansion_appendix}
\end{equation}

The leading coefficient $2$ is the continuum expression of the two serial endpoint constraints required to complete one nearest-neighbor transfer.

Internal phase mismatch may alter excitation amplitude, polarization, mode participation, or defect coupling. It does not supply an additional universal multiplicative factor in the coherent transverse propagation speed unless such a factor is separately derived from the microscopic eigenvalue problem.

\section{Saturation and Loss of the Active Coordinate}
\label{sec:grav_saturation_appendix}

In the unstrained active vacuum,
\begin{equation}
    \alpha_{\mathrm{G}} = 0,
\end{equation}
so
\begin{equation}
    \mathcal{C} = 1, \qquad \mu_{\mathrm{eff}} = \mu_{\mathrm{D}}, \qquad c_{\mathrm{P}} = c_{\mathrm{E}}.
\end{equation}

As the active FCC phase approaches saturation,
\begin{equation}
    \alpha_{\mathrm{G}} \rightarrow 1^{-},
\end{equation}
one has
\begin{equation}
    \mathcal{C} \rightarrow 0^{+}.
\end{equation}

The referred surface inertia diverges:
\begin{equation}
    \mu_{\mathrm{eff}} = \frac{\mu_{\mathrm{D}}}{\mathcal{C}^2} \rightarrow \infty,
\end{equation}
while the completed transverse propagation speed vanishes:
\begin{equation}
    c_{\mathrm{P}} = c_{\mathrm{E}}\mathcal{C} \rightarrow 0.
\end{equation}

The mechanical interpretation is
\begin{equation}
    \text{finite active-side tangent stiffness} + \text{divergent referred inertia} = \text{causal arrest}.
\end{equation}

The divergence does not represent an ordinary infinite material density of the SC substrate. It indicates that the completed active-FCC coordinate becomes singular and loses validity at the phase boundary.

The present construction applies only while both endpoints remain active:
\begin{equation}
    0 \leq \alpha_{\mathrm{G}}^{(i)} < 1, \qquad 0 \leq \alpha_{\mathrm{G}}^{(j)} < 1.
\end{equation}

The FCC endpoint constraint is not analytically continued through the FCC-to-SC phase transition.

\section{Microscopic Completion Criterion}
\label{sec:grav_microscopic_completion}

The present appendix gives a conservative reduced derivation of the gravitational endpoint constraint and its kinetic and dispersive consequences. It does not complete the underlying parent microphysics.

A complete node--contact--condensate--Tension-Column action must derive, rather than impose, the single-endpoint relation
\begin{equation}
    u_{i} = \left( 1-\alpha_{\mathrm{G}}^{(i)} \right)y_{i}.
\end{equation}

After elimination of any auxiliary condensate and vertical-constraint variables, the resulting reduced primitive-cell dynamics must reproduce or revise the transfer operator
\begin{equation}
    \mathbf{T}_{\mathrm{cell}} = \mathbf{I} + \left( \mathcal{C}-1 \right) \mathbb{P}_{\mathrm{tr}},
\end{equation}
and the referred mass matrix
\begin{equation}
    \mathbf{M}_{\mathrm{cell}}^{(\mathrm{eff})} = \mathbf{T}_{\mathrm{cell}}^{-\mathsf{T}} \mathbf{M}_{\mathrm{cell}}^{(0)} \mathbf{T}_{\mathrm{cell}}^{-1}.
\end{equation}

The long-wavelength coherent transverse branch must satisfy
\begin{equation}
    \lim_{k\rightarrow0} \frac{ \omega_{\mathrm{T}} \left( k, \alpha_{\mathrm{G}} \right) }{k} = c_{\mathrm{E}} \left( 1-\alpha_{\mathrm{G}} \right)^2.
    \label{eq:microscopic_completion_speed}
\end{equation}

Equivalently,
\begin{equation}
    \frac{ c_{\mathrm{P}}^2 \left( \alpha_{\mathrm{G}} \right) }{ c_{\mathrm{E}}^2 } = \left( 1-\alpha_{\mathrm{G}} \right)^4.
    \label{eq:microscopic_completion_ratio}
\end{equation}

The complete finite-wave-number theory must also verify
\begin{equation}
    \omega_{s}^2 \left( \mathbf{k}, \alpha_{\mathrm{G}} \right) \geq 0
\end{equation}
for every physical branch $s$, every wave vector in the first Brillouin zone, and every background satisfying
\begin{equation}
    0 \leq \alpha_{\mathrm{G}} < 1.
\end{equation}

If the completed parent action produces a different endpoint dependence, transfer operator, inertia scaling, or long-wavelength dispersion, then the reduced closure developed here must be modified or rejected.

\section{Conclusion}

\hyperref[app:fcc_microphysics]{Appendix~\ref*{app:fcc_microphysics}} establishes the gravity-independent FCC contact carrier. The present appendix applies the gravitational endpoint constraint to that carrier and derives the resulting kinetic and dispersive reduction.

The completed mechanical chain is
\begin{equation}
\begin{array}{c}
    \tau_{\mathrm{av}}^{(i)} = \tau_{\mathrm{D}} \left( 1-\alpha_{\mathrm{G}}^{(i)} \right) \\
    \Downarrow \\
    u_{i} = \left( 1-\alpha_{\mathrm{G}}^{(i)} \right)y_{i} \\
    \Downarrow \\
    \mathcal{C}_{ij} = \left( 1-\alpha_{\mathrm{G}}^{(i)} \right) \left( 1-\alpha_{\mathrm{G}}^{(j)} \right) \\
    \Downarrow \\
    \mathbf{M}_{\mathrm{cell}}^{(\mathrm{eff})} = \mathbf{T}_{\mathrm{cell}}^{-\mathsf{T}} \mathbf{M}_{\mathrm{cell}}^{(0)} \mathbf{T}_{\mathrm{cell}}^{-1} \\
    \Downarrow \\
    \rho_{\mathrm{eff}} = \dfrac{ \rho_{\mathrm{eff}}^{(0)} }{ \mathcal{C}^2 } \\
    \Downarrow \\
    \omega_{\mathrm{T}}^2 = c_{\mathrm{E}}^2 \mathcal{C}^2k^2 + \mathcal{O} \left( k^4d_{\mathrm{nn}}^2 \right) \\
    \Downarrow \\
    c_{\mathrm{P}} = c_{\mathrm{E}} \left( 1-\alpha_{\mathrm{G}} \right)^2.
\end{array}
\end{equation}

At the primitive-cell microscopic scale,
\begin{equation}
    \rho_{\mathrm{eff}}^{(0)} = \rho_{\mathrm{sh}}^{(\mathrm{micro},0)}, \qquad K_{\mathrm{sh}}^{(0)} = K_{\mathrm{sh}}^{(\mathrm{micro},0)},
\end{equation}
while $\rho_{\mathrm{eff}}^{(0)}$ and $K_{\mathrm{sh}}^{(0)}$ remain the scale-neutral notation when the same constrained-coordinate structure is written at a more general continuum level.

The endpoint law is therefore not left as a bare propagation postulate. It is represented as a reduced work-conjugate coordinate constraint whose referred inertia and long-wavelength dispersion follow from a conservative constrained action. The remaining open problem is the derivation of that constraint from the complete microscopic parent dynamics.
